\section{Minimal Core Language}
\label{sec:core}

\subsection{Primitive objects}

The language has four primitive objects:

\begin{definition}[Core language]
$(\mathrm{At}, F, \mathcal{B}, \mathcal{A})$ where:
\begin{itemize}
  \item $\mathrm{At}$: atomic propositions (world-states).
  \item $F$: rule-based generation operator (closure / semantics).
  \item $\mathcal{B}$: break family (attack surface; source of invalidation).
  \item $\mathcal{A}$: adoption dynamics (how agents select interpretations).
\end{itemize}
\end{definition}

\begin{definition}[Canonical world]
$P\subseteq\mathrm{At}$ is canonical ($\mathrm{SC}_T(P)=1$) if $F(P)\subseteq P$
and $P$ has no alternative admissible fixed point.
\end{definition}

\begin{definition}[Break cost]
$C_{\min}(P) = \min\text{-HS-cost}(\mathcal{B}(P))$.
If $\mathcal{B}(P)=\emptyset$: $C_{\min}(P)=+\infty$.
\end{definition}

\subsection{Justification structure}

A justification $J\in\mathcal{J}_{\mathrm{valid}}(P)$ is a certificate that
atom $a\in P$ is derivable. The Break set $\mathrm{Break}(J)$ is the
minimal set of source nodes whose disabling invalidates $J$.

\begin{definition}[Unavoidable external seed]\label{def:ues}
$a\in\mathrm{UES}(P)$ if every justification $J$ for $a$ has
$\mathrm{Break}(J)\neq\emptyset$ (cannot be fully internally certified).
\end{definition}
